Clinical adoption of artificial intelligence requires calibrated reliance rather than blind
trust or a persuasive explanation attached to every output. Clinicians should be able to
judge when an AI output has a defensible basis, when it should be challenged, and where the
model's competence ends. Explanations that fit the decision context can support rational
checking \cite{tonekaboni2019}, but familiar explanations can also increase reliance when AI
advice is wrong \cite{prinster2024}. A clinically recognizable explanation is therefore not
evidence of trustworthy model behavior unless the proposed link can itself be tested.

This is the motivating problem for prostate-cancer AI: beyond asking whether a model can
separate clinically relevant states, we need to determine which human-interpretable
quantitative axes it represents, whether each axis remains recoverable across cohorts and
technical variation, and whether a downstream decision actually uses it. Answering these
questions axis by axis can make model reliability inspectable in terms clinicians can contest,
rather than reducing trustworthiness to predictive performance or a plausible visualization.

Clinically established quantitative and ordinal targets offer a possible interface for such
testing. A pathologist can define, measure, contest, and communicate grade or tumor extent;
molecular laboratories can provide assay-defined targets; and clinical records provide
outcomes. If an AI signal is demonstrably linked to one of these targets, the proposed basis
becomes inspectable: a pathologist can confirm that the tissue supports the named feature,
reject the explanation when it does not, or identify a meaningful discordance. Concept-based
methods provide precedents. TCAV tests prediction sensitivity to human-defined concepts, and
concept bottleneck models make named concepts explicit and intervenable
\cite{kim2018,koh2020}. Medical-imaging and pathology systems have likewise exposed
confounders or connected influential regions and embeddings to expert-recognizable findings
\cite{sauter2022,lou2026,han2026}. These studies also show why concept labels, embedding
similarity, and actual decision dependence cannot be treated as interchangeable.

Three evidence stages are particularly easy to collapse. A target may be recoverable from a
frozen representation; that relationship may reproduce across cohorts and technical
settings; and a downstream disease head may functionally use the target. Only the third
directly tests sensitivity of a particular prediction rule, and even faithful functional use
would not prove clinical benefit. We therefore order the evidence as recoverability,
reproducibility and qualification, functional-use testing, and finally prospective evaluation
of clinician reliance and utility. Each stage enables, but does not replace, the next.

\paragraph{Clinical orientation and reference availability.}
The disease studied here is predominantly prostate adenocarcinoma. Routine hematoxylin-and-
eosin (H\&E) sections expose gland architecture, tumor amount, and tissue context but do not
directly measure DNA alterations or transcriptional activity. The six study axes therefore
span different biological levels \cite{panda2022,tcgaprad2015,tcgacdr2018}. Grade/ISUP measures
architectural differentiation; tumor phenotype/content measures tumor presence or
amount; PTEN loss and SPOP mutation are binary molecular alterations; AR activity is a
continuous transcriptional pathway score. Recurrence is a patient outcome observed after
treatment rather than a microscopic structure graded on the slide. They are not six
interchangeable measures of ``cancer severity.''

Crucially, the slides were model inputs, not the only study data. Every reported alignment
estimate paired a slide-derived frozen embedding with a target-specific reference at the same
analysis unit. The encoder was not fine-tuned on these study labels; a lightweight probe used
the paired references to test whether the embedding contained recoverable target
information. Reference availability nevertheless differed: morphology had compatible
external labels, molecular references were largely confined to TCGA-PRAD, and recurrence
definitions were non-equivalent across resources. Missing reference data were treated as not
evaluable, not as evidence of non-alignment. Unsupported required an available comparison
that did not qualify the proposed signal, as for SPOP. The detailed target, reference,
availability, metric, and cohort frames are provided in Supplementary Section S1.

\paragraph{Why the six prostate-cancer axes are not interchangeable.}
A positive result supports a different claim at each level. Grade and phenotype are direct
morphology targets; PTEN, SPOP, and AR are indirect histology--molecular associations; and
recurrence is shaped by biology, treatment, follow-up, censoring, and endpoint definition.
Recovering a target shows that related information is accessible in the embedding, not that
H\&E replaces a molecular assay or that a later prediction head uses the target. Conversely,
failure to recover or transport a target warns against invoking it as an explanation, but
does not prove absence of the underlying biology.

Existing studies often summarize image--label association with one pooled correlation or
area under the receiver operating characteristic curve. Such a number cannot distinguish
target information from grade composition, institutional case mix, acquisition conditions,
sampling, or endpoint definitions \cite{dawood2026,neidlinger2025,pathorob2026}. Nor does it
show whether a fitted relationship survives a new cohort without refitting. We therefore
define \emph{representation--target alignment} as patient-level association or
discrimination between a registered target and the out-of-fold or locked external output of
a simple probe fitted to a frozen CONCH or Virchow representation
\cite{conch2023,virchow2023}. We qualify each target across recoverability, external transport,
conditional increment, site behavior, representation and sampling sensitivity, and endpoint
fidelity, retaining these as an evidence vector rather than a composite score.

\paragraph{Position of this study.}
This study compares Grade/ISUP, tumor phenotype/content, PTEN loss, SPOP mutation, AR
activity, and recurrence to determine which targets are accessible in frozen pathology
representations and how their evidence changes across cohorts, clinical adjustment, sites,
representations, sampling, and endpoint definitions. Grade and phenotype provide the
broadest test of external transport; PTEN, AR, and SPOP distinguish recoverable, conditional,
and unsupported molecular states; and recurrence tests whether outcome-linked information
retains its meaning across event definitions and comparators. Among these axes, Grade/ISUP
has the most complete evidence chain because it is additionally tested in two locked
biochemical-recurrence heads and in site-heldout and independent-patient functional-transport
analyses. Internal sensitivity was not promoted because site evidence was encoder-specific and
the external gate failed or remained inconclusive. ISUP is thus the most developed functional example, not the
organizing target. The study does not explain a deployed model, establish indispensable
or externally replicated functional use, comprehensively test clinical increment, or measure
clinician reliance or benefit.

Its primary task is comparative: to map the evidential status of all six axes without treating
their labels, biological levels, or available validation opportunities as equivalent.

\paragraph{Contributions.}
This paper makes four contributions. First, it defines clinically interpretable targets as
\emph{contestable shared coordinates} through which a pathologist can agree or disagree with
a clinically stated basis for an AI output. Second, it provides a comparative six-axis map
across cohorts and two foundation models, distinguishing transported morphology,
context-sensitive molecular and outcome associations, and unsupported evidence. Third, it
turns the map into an actionable audit vocabulary: transported coordinates can anchor
clinician review of agreement and disagreement; conditional or unsupported axes indicate
where explanations should be withheld; and evidence gaps prioritize external or functional
validation. Fourth, it provides a worked extension from representation availability toward
functional testing by applying targeted erasure and matched controls to an
ISUP-correlated direction in locked CONCH- and Virchow-derived risk heads, then testing its
site and independent-patient transport without target-cohort refitting. This final
experiment extends rather than organizes the six-axis comparison. The proposed audit
vocabulary is intended for model evaluation, communication, and validation planning; whether
showing it to clinicians improves decisions remains untested. Establishing this known-target
map also defines the prerequisite for residual discovery: only after reliable shared coordinates
and technical confounders are accounted for should remaining reproducible, decision-linked
signals be investigated as candidate biomarkers rather than artifacts.
